%% This is an example first chapter.  You should put chapter/appendix that you
%% write into a separate file, and add a line \include{yourfilename} to
%% main.tex, where `yourfilename.tex' is the name of the chapter/appendix file.
%% You can process specific files by typing their names in at the 
%% \files=
%% prompt when you run the file main.tex through LaTeX.
\chapter*{List of articles and publications}

This thesis focuses on the development of new approaches based on adaptive control in the presence of sensor and actuator faults for UAVs. The thesis was conducted under joint supervision with the LAMIH – UMR CNRS 8201, in collaboration with Dr. Chouki Sentouh.\\
Within the framework of this collaboration, I gained extensive knowledge in the field of UAVs. Furthermore, several approaches were developed and experimentally validated on UAV models. The results obtained were submitted to and accepted by several conferences and scientific journals.
\subsection*{Published Papers in an International Journal}
\begin{itemize}
    \item[1.] \textbf{Bacha, A}., Chelihi, A., Glida, H. E., Sentouh, C., $\&$ Senoussi, M. A. (2025). Optimal fault-tolerant adaptive fuzzy control of quadrotor UAV: a fixed-time stability approach. \textit{International Journal of Dynamics and Control}, 13(5), 1--16, [IF=1.9].
    \item[2.] \textbf{Bacha, A}., Chelihi, A., Glida, H. E., $\&$ Sentouh, C. (2024). Fixed-time fault-tolerant adaptive neural network control for a twin-rotor UAV system with sensor faults and disturbances. \textit{Drones}, 8(9), 467, [IF=4.8].
    % \textbf{Fixed-time adaptive neural tracking control for a helicopter-like twin rotor MIMO system}
    % \begin{itemize}
    %     \item[] \textit{Conference:} The Third International Nonlinear Dynamics Conference (NODYCON 2023) 01 May 2024
    %     \item[] \textit{DOI:} \url{https://doi.org/10.1007/978-3-031-50635-2_16}
    % \end{itemize}
    
    % \item \textbf{Fixed-time fault-tolerant adaptive neural network control for a twin-rotor uav system with sensor faults and disturbances}
    % \begin{itemize}
    %     \item[] \textit{Journal:} Drones 8 September 2024 \textbf{(MDPI: A Class)}
    %     \item[] \textit{DOI:} \url{https://doi.org/10.3390/drones8090467}
    % \end{itemize}
    
    % \item \textbf{Adaptive Fuzzy Control for a Quadrotor UAV Attitude with Actuator and Sensor Failure The Practical Fixed-Time Stability}
    % \begin{itemize}
    %     \item[] \textit{Conference:} International Conference on Control, Decision and Information Technologies (CoDIT) 18 October 2024
    %     \item[] \textit{DOI:} \url{https://doi.org/10.1109/CoDIT62066.2024.10708478}
    \end{itemize}
    \subsection*{International Conference}
    \begin{itemize}
    \item[1.]  BACHA, Aymene, CHELIHI, Abdelghani, GLIDA, Hossam Eddine, et SENTOUH, Chouki. Adaptive fuzzy control for a quadrotor uav attitude with actuator and sensor failure: the practical fixed-time stability. In : 2024 10th International Conference on Control, Decision and Information Technologies (CoDIT). IEEE, 2024. p. 504-509.
    \item[2.] BACHA, Aymene, CHELIHI, Abdelghani, et SENTOUH, Chouki. Fixed-Time Adaptive Neural Tracking Control for. In : Advances in Nonlinear Dynamics, Volume III: Proceedings of the Third International Nonlinear Dynamics Conference (NODYCON 2023). Springer Nature, 2024. p. 163. 
    % \item \textbf{Optimal fault-tolerant adaptive fuzzy control of quadrotor UAV a fixed-time stability approach}
    % \begin{itemize}
    %     \item[] \textit{Journal:} International Journal of Dynamics and Control 13 July 2023 \textbf{(Springer: A Class)}
    %     \item[] \textit{DOI:} \url{https://doi/abs/10.1002/acs.3659}
    % \end{itemize}
\end{itemize}
\subsection*{Submitted and Under Review}
\begin{itemize}
    \item[1.] Senoussi. M.A, Boumahrez. M, Vicenç. P, Glida. H.E $\&$ \textbf{Bacha, A}. Data-driven fault-tolerant MPC for quadrotors: a quasi linear parameter varying approach with RBF neural network compensation. % \textit{Drones}, 8(9), 467, [IF=4.8]. 
    % \item[2.]
\end{itemize}
\chapter{General Introduction}
\label{chap:introduction}

% --- Section 1.1: The Rise of Unmanned Aerial Vehicles ---
\section{General Context and Motivation}
\label{sec:general_context}

Over the past two decades, the field of autonomous systems has witnessed a significant paradigm shift. The rapid growth in the use of Unmanned Aerial Vehicles (UAVs) represents one of the most prominent manifestations of this transformation \cite{UAV_rev,UAV_rev1}. Once confined primarily to military applications, UAVs have evolved into indispensable tools across a wide range of civilian, scientific and commercial domains. This evolution has been driven by parallel advances in micro-electro-mechanical systems, embedded computing platforms, battery technologies and intelligent control algorithms. As a result, modern UAVs exhibit remarkable operational flexibility, autonomy and efficiency.

The societal and industrial impact of UAVs is substantial, encompassing numerous sectors, as summarized in Table~\ref{tab:uav_applications}. In precision agriculture, UAVs provide high-resolution data for crop monitoring, irrigation optimization and targeted spraying, thereby improving productivity while reducing environmental impact \cite{agrcltr1,agrcltr2,agrcltr3}. In civil engineering, they are widely employed for infrastructure inspection tasks in hazardous or inaccessible environments, enhancing safety and significantly lowering operational costs \cite{UAVs_civil_eng1,UAVs_civil_eng2,UAVs_civil_eng3}. Moreover, UAVs play a critical role in public safety operations, particularly during disaster response missions such as search and rescue, damage assessment and emergency supply delivery \cite{Emergency_response1,Emergency_response2}. These applications, among others, illustrate the transformative influence of UAV technologies across modern society.


\begin{table}
\centering
\caption{Key Fields of Application for UAVs}
\label{tab:uav_applications}
\renewcommand{\arraystretch}{1.05} % reduce row spacing
\setlength{\tabcolsep}{6pt} % reduce column spacing
\begin{tabular}{|p{0.25\textwidth}|p{0.65\textwidth}|}
\hline
\textbf{Field of Use} & \textbf{Applications} \\
\hline
Agriculture \& Environment & 
Crop health, precision spraying, wildlife tracking \\
\hline
Civil Infrastructure & 
Inspection of bridges, power lines, site surveying, pipeline and turbine monitoring \\
\hline
Public Safety \& Disaster Mgmt. & 
Search/rescue, post-disaster assessment, wildfire detection, medical supply delivery \\
\hline
Commercial \& Media & 
Package delivery, aerial cinematography, real estate tours, live broadcasting \\
\hline
Scientific Research & 
Atmospheric data, geological mapping, oceanographic studies \\
\hline
\end{tabular}
\end{table}

Among the various UAV configurations, rotorcraft platforms have attracted particular attention due to their Vertical Takeoff and Landing (VTOL) capability. This feature enables hovering and operation in confined environments, making rotorcraft UAVs highly suitable for urban, industrial and indoor missions. Nevertheless, the same characteristics that grant such versatility also introduce significant control challenges. Rotorcraft UAVs are inherently nonlinear, strongly coupled, and often unstable, which necessitates the development of advanced, robust and adaptive control strategies to ensure safe and reliable operation.


\section{Core Control Challenges}
\label{sec:intro_challenge}

Despite their structural differences, the rotorcraft platforms under investigation share a common set of control challenges:
\begin{itemize}
    \item \textbf{Nonlinear and Coupled Dynamics \cite{NonlinearCoupledDynamics1,NonlinearCoupledDynamics2}:} The relationship between rotor thrust and the vehicle's motion is highly nonlinear. Furthermore, the translational and rotational dynamics are strongly coupled, meaning that a change in one state invariably affects others.
    \item \textbf{Underactuation \cite{Underactuation1,Underactuation2}:} Most rotorcraft, including the quadrotor, are underactuated systems. They have more degrees of freedom (e.g., 6-DOF for position and orientation) than independent control inputs (e.g., four rotors). This necessitates indirect control strategies and complex state management.
    \item \textbf{System Uncertainties and External Disturbances \cite{ExternalDisturbances1,ExternalDisturbances2}:} UAVs operate in unpredictable environments. Their performance is constantly affected by external disturbances such as wind gusts, as well as internal uncertainties like variations in payload, battery discharge, and unmodeled aerodynamic effects.
\end{itemize}

Beyond these environmental factors, the UAV faces a more specific internal risk: hardware faults in the actuators and sensors. These components can malfunction during flight, effectively changing the physical behavior of the vehicle. This creates a dual problem where the system must handle not only unpredictable wind and payload changes but also a vehicle that is reacting differently than expected. This combination presents a significant control challenge, as the UAV must remain stable even when its own hardware is no longer behaving as it was originally designed.

% --- Section 1.4: The Critical Need for Fault Tolerance ---
\section{The Need for Fault-Tolerant Control}
\label{sec:intro_ftc}

The successful use of autonomous UAVs depends on their ability to operate safely and reliably. The biggest threat to this reliability is a fault in the system's hardware, especially in its sensors and actuators \cite{fault}. A fault, whether it is a sudden failure or a gradual degradation over time, can have severe consequences, possibly leading to poor performance, instability, or a complete loss of the vehicle.

Actuator faults, such as a motor losing some of its power or a propeller malfunction, can create an imbalance of forces and torques that the controller must quickly correct \cite{A_fault}. Sensor faults, like incorrect readings (bias), gradual errors over time (drift), or random noise from an Inertial Measurement Unit (IMU) or GPS, can give the control system wrong information \cite{S_fault}. This can cause the controller to make incorrect and dangerous decisions.

For these reasons, creating Fault-Tolerant Control (FTC) systems is not just an optional improvement but a basic need for safe and reliable UAV operations \cite{ftc_UAV_review1,ftc_UAV_review2}. An effective FTC system must be able to find, identify, and correct for faults as they occur. This ensures that the UAV can stay stable and finish its mission, or at the very least, perform a safe emergency landing.

However, simply detecting and compensating for a fault is not enough; the speed of this response is a critical factor that is often overlooked. When a component fails during flight, the UAV’s stability can deteriorate rapidly, leaving only a very narrow time window to intervene before the vehicle becomes unrecoverable. Standard control methods, which typically guarantee convergence merely over an infinite horizon or depend heavily on initial conditions, may be too slow to prevent a crash in these high-stakes scenarios. Consequently, to ensure safety, the system requires the property of fixed-time stability. This theoretical guarantee ensures that control errors converge to zero within a pre-calculated, bounded time regardless of the magnitude of the initial fault, providing a deterministic upper bound for the system's recovery.

% --- Section 1.5: Research Objectives and Contributions ---
\section{Research Objectives and Contributions}
\label{sec:intro_objectives}

The primary objective of this thesis is to design and validate advanced, robust, and fault-tolerant control strategies for rotorcraft UAVs that can ensure high-performance trajectory tracking in the presence of unknown dynamics, external disturbances, and critical faults in both sensors and actuators.

\noindent To achieve this goal, this thesis makes the following key contributions:

\begin{enumerate}
    \item \textbf{Development of Fault-Tolerant Fixed-Time Stable Controllers:} A core contribution of this work is the application of fixed-time stability theory. Unlike traditional controllers that guarantee stability over an infinite or state-dependent time horizon, our proposed controllers ensure that the UAV converges to its desired trajectory and recovers from disturbances or faults within a predictable and bounded time, regardless of the initial conditions \cite{Polyakov2012}. This provides a strong, verifiable guarantee of performance, which is essential for safety-critical missions \cite{Ni2016}.

    \item \textbf{Intelligent Model-Free Control Approaches:} This thesis moves beyond traditional model-based control by leveraging intelligent function approximators. We employ adaptive Radial Basis Function Neural Networks (RBFNN) and Fuzzy Logic Systems (FLS) to estimate and compensate for the complex, unknown nonlinearities and fault dynamics in real-time. This model-free approach enhances the controller's adaptability and robustness, relying on the universal approximation capabilities of these structures rather than a precise mathematical model \cite{Park1991, Wang1994}.

    \item \textbf{Optimization for Enhanced Performance:} The performance of advanced controllers is highly dependent on the tuning of their parameters. This thesis integrates metaheuristic optimization algorithms, specifically the Grey-Wolf Optimizer (GWO) and the Whale Optimization Algorithm (WOA), to automatically and systematically tune the controller gains. This optimization process ensures the best possible tracking accuracy and control smoothness, overcoming the limitations of manual tuning methods often cited in control literature \cite{Precup2015, Eltogby2016}.
\end{enumerate}% --- Section 1.6: Thesis Outline ---
\section{Thesis Outline}
\label{sec:intro_outline}

The remainder of this thesis is structured as follows:

\begin{itemize}
    \item \textbf{Chapter 2} provides a comprehensive review of the state-of-the-art in UAVs control and fault tolerance. It also introduces the key theoretical preliminaries, including nonlinear control techniques, intelligent approximators, the formal definition of fixed-time stability and the metaheuristic optimization algorithms.
    \item \textbf{Chapter 3} presents the design and validation of an active fault-tolerant fixed-time adaptive neural network controller for the Twin-Rotor MIMO System (TRMS). This chapter focuses on addressing sensor faults in a highly coupled, helicopter-like system.
    \item \textbf{Chapter 4} details the development of an optimal fault-tolerant adaptive fuzzy control scheme for a quadrotor UAV. This work addresses the challenge of handling simultaneous faults in both sensors and actuators.
    \item \textbf{Chapter 5}  extends the fault-tolerant control framework to a multi-UAV formation. It addresses the leader-follower formation problem, proposing a distributed fixed-time control protocol that ensures all follower UAVs maintain a desired geometric formation relative to the leader. The design incorporates a consensus-based estimator to handle communication delays and an adaptive mechanism to compensate for actuator faults within the formation, ensuring collision-free and synchronized flight. 
        
       
\end{itemize}

%\end{document}

